% CS 418: Intro to Data Science — Worksheet 6 (06-week03-day2)
% Week 3, Wednesday 2026-09-09
% Compile with: latexmk -pdf worksheet.tex

\documentclass[11pt]{article}

% ── Packages ──────────────────────────────────────────────────────────────────
\usepackage[margin=0.65in, headheight=14pt]{geometry}
\usepackage{amsmath, amssymb}
\usepackage{ragged2e}
\usepackage{enumitem}
\usepackage{booktabs}
\usepackage{graphicx}
\usepackage{hyperref}
\usepackage{xcolor}
\usepackage{fancyhdr}
\usepackage{mdframed}
\usepackage{array}

% ── Page style ────────────────────────────────────────────────────────────────
\pagestyle{fancy}
\fancyhf{}
\lhead{\textbf{CS 418: Intro to Data Science}}
\rhead{UIC, Fall 2026}
\cfoot{\thepage}
\renewcommand{\headrulewidth}{1pt}
\setlength{\headheight}{12pt}
\setlength{\headsep}{0.4cm}

% ── Hyperlinks ────────────────────────────────────────────────────────────────
\hypersetup{colorlinks=false, hidelinks}

% ── Convenience environments ──────────────────────────────────────────────────
\newenvironment{problem}[1]{%
  \begin{mdframed}[backgroundcolor=white, linecolor=black, linewidth=1pt]
  \textbf{Problem #1.}
}{%
  \end{mdframed}\vspace{0.5em}
}

\newenvironment{hint}{%
  \begin{mdframed}[backgroundcolor=white, linecolor=gray, linewidth=0.5pt]
  \textit{Hint:}
}{%
  \end{mdframed}\vspace{0.3em}
}

% Usage: \blankspace{6cm} — blank area for hand-written work, no rules
\newcommand{\blankspace}[1]{%
  \vspace{0.3em}
  \begin{minipage}[t][#1][t]{\linewidth}
  \end{minipage}
  \vspace{0.3em}
}

% Usage: \guidedopen{title}{guided content}{guided workspace height}{open-ended workspace height}
% Left: a titled, structured prompt with its own blank workspace. Right: a
% fully blank workspace for open-ended exploration of the same topic (no
% prompt text), sized so the two columns land at roughly the same total
% height. Neither workspace is ruled.
\newcommand{\guidedopen}[4]{%
  \noindent\textbf{#1}\\[0.18em]
  \noindent
  \begin{minipage}[t]{0.485\textwidth}
    #2
    \blankspace{#3}
  \end{minipage}
  \hfill
  \begin{minipage}[t]{0.485\textwidth}
    \blankspace{#4}
  \end{minipage}
  \\[0.4em]
}

% ══════════════════════════════════════════════════════════════════════════════
\begin{document}
\RaggedRight

% ── Title block: topic left, info box right ───────────────────────────────────
\noindent
\begin{minipage}[t]{0.55\textwidth}
  \vspace{0pt}
  {\LARGE \textbf{Worksheet \#\,6}} \\[0.18em]
  {\large \textit{Week 3, Day 2: Wrangling, data formats}}
\end{minipage}
\hfill
\begin{minipage}[t]{0.40\textwidth}
  \vspace{0pt}
  \small
  Name: \underline{\hspace{4.5cm}} \\[0.35em]
  NetID: \underline{\hspace{4.3cm}} \\[0.35em]
  Date: \underline{\hspace{4.4cm}} \\
\end{minipage}

\vspace{0.6em}
\noindent\rule{\linewidth}{0.6pt}
\vspace{0.1em}

\noindent\underline{\textbf{Guided checkpoints (optional)}}
\vspace{0.8em}

% ── 1. Types of quality issues ──────────────────────────────────────────────
\guidedopen{1. Types of Quality Issues}{
List the  types of data quality issues
}{3.2cm}{4.0cm}

\vfill

% ── 2. The granularity spectrum ─────────────────────────────────────────────
\guidedopen{2. The Granularity Scale}{
Draw the scale from fine-grained to coarse-grained data, with examples.
}{4.0cm}{4.8cm}

\vfill

% ── 3. Group by, pictured ───────────────────────────────────────────────────
\guidedopen{3. Group By, Visualized}{
Sketch the three steps of ``group by.''
}{4.3cm}{5.2cm}

\newpage

% ── 4. Pivot: one cell per combination ──────────────────────────────────────
\guidedopen{4. Pivot: One Cell per Combination}{
Sketch the pivoted grid (one cell comes up empty)
}{4.0cm}{4.8cm}

\vfill

% ── 5. Is the data in scope? ────────────────────────────────────────────────
\guidedopen{5. Is the Data in Scope?}{
Draw the two scope visuals.
}{4.0cm}{5.0cm}

\vfill

% ── 6. File formats ─────────────────────────────────────────────────────────
\guidedopen{6. File Formats}{
List the data formats discussed.
}{3.2cm}{4.0cm}

\vfill

\end{document}